\documentclass{beamer}

\usepackage{pifont}
\usepackage{tikz}
\usetikzlibrary{arrows.meta,positioning, calc}
\usepackage{caption}
\usepackage[ngerman]{babel}
\usepackage{csquotes}

\newcommand{\ideaicon}{\ding{72}}
\newcommand{\screenicon}{\ding{229}}
\newcommand{\personicon}{\ding{110}}
\newcommand{\checkicon}{\ding{51}}
\newcommand{\warnicon}{\ding{115}}

\usetheme{metropolis}
\setbeamertemplate{footline}{}
\title{Offene Daten erlebbar machen}
\date{}
\author{Adaktylos Leonidas Odysseas, Baumgartner Vincent Vitez,\\ Neda Nikkhousarighiyeh, Alexander Grath}

\begin{document}
  %\maketitle

  \begin{frame}{Benutzertests – Methode}
    \textbf{8 Testpersonen}, je 2 pro Teammitglied \quad | \quad ca. 30–45 Min. pro Sitzung

    \vspace{0.4cm}
    \begin{tabular}{@{}ll@{}}
      \textbf{Methode} & Think-Aloud, Post-Task-Fragen, SUS, Interview \\[0.4em]
      \textbf{Aufgabe 1} & Kunsthistorisches Museum auf der Karte finden \\
      \textbf{Aufgabe 2} & Datensatz \emph{Radwege Wien} herunterladen \\
      \textbf{Aufgabe 3} & Schloss Schönbrunn über die Suche finden \\
      \textbf{Aufgabe 4} & Dunkelmodus aktivieren + Datensatz teilen \\
    \end{tabular}
  \end{frame}

  \begin{frame}{Motivation}
    \begin{itemize}
      \setlength\itemsep{0.6em}
      \item Wien veröffentlicht \textbf{über 300 offene Datensätze} auf dem städtischen Open-Data-Portal
      \item Nur \textbf{8 mobile Apps} visualisieren diese Daten – meist veraltet oder spezialisiert
      \item Problem: Daten sind verfügbar, aber \textbf{nicht zugänglich}
      \item Lösung: \textbf{Dat-O-Kart} – farbkodierte Pins, Offline-Download, Suche, Dunkelmodus
    \end{itemize}
  \end{frame}

  \begin{frame}{Related Work}
    \begin{itemize}
      \setlength\itemsep{0.5em}
      \item Eine Studie zu Open-Government-Data zeigt: die meisten Portale bieten keine visuelle Aufbereitung
      \item Gut gestaltete Dashboards erhöhen Verständlichkeit – heterogene Nutzergruppen berücksichtigen
      \item Farbige Symbole bleiben besser in Erinnerung → farbkodierte Pins
      \item \textbf{Konkurrenz}: Wien Karte, Eurostat, OECD Explorer – keine mobile Open-Data-Exploration für Wien
    \end{itemize}
  \end{frame}


\end{document}
